\section{Consumer Packaging and Release Diagnostics}
\label{sec:consumer_packaging}
\subsection{Packaging recipe and importer check}
The upstream repository root can differ from the directory selected for packaging. We examined the catalog-pinned YunoHost recipes for 15 unresolved projects with empty pnpm workspace roots and 16 without a conventional root lockfile. We checked source declarations, install and upgrade instructions, and available service templates. Twelve of the 15 workspace recipes install a release artifact or distribution package. Of the remaining three, one delegates asset compilation without selecting a named importer in the observed recipe, one builds a named frontend workspace, and one builds a workspace and copies a named client into a named server. Three of the 16 missing-root-lockfile recipes contain explicit npm dependency commands. One uses a nested server lockfile. The other two lack a committed npm lockfile at the recipe-selected root.

We resolved the source references for all six recipes with explicit JavaScript package commands. Four select commits different from the upstream commits observed in the root study. The three recipes with named importers and lockfiles yield four importer paths across three applications. We verified each declared source archive checksum and checked that the selected manifest and lockfile bytes match the pinned Git files. These three packaging trees contain no source patches in the checked locations. The comparison traces mandatory dependency declarations from the selected importers, using the same three portfolios and tie policy. It describes the packaging recipe's inputs. Runtime inclusion and installed versions would require separate evidence.

We corroborated 425 paths across the four importers and obtained usable maintainer lists for all 401 requested versions. All three applications had documented hits under both portfolios for every version cutoff choice. In the same-commit case, the recipe-selected frontend supplies paths that the empty upstream root cannot provide. The other two recipes use different commits from the root study. This diagnostic uses a separate three-application denominator. The 215-project comparison retains its root-importer scope and bounds.

\subsection{Same-commit comparison on a fresh recipe sample}
We froze 99 previously unaudited recipes from the catalog frame. The sample includes all 39 remaining repositories with an observed workspace marker and a seeded draw of 60 from the other 430. Seven recipes name a non-root importer involved in a build or service (six in the workspace stratum and one in the other stratum). Six of the seven select Git commits different from the screened upstream HEADs. Four source archives and their corresponding Git files could be checked by hash. One of the four has packaging patches and another uses an unsupported Yarn workspace. The remaining two permit a comparison at a fixed recipe-selected commit, using the same 772 exact-version requests across both views of the application. We obtained 771 usable maintainer lists.

In one application, the root yielded 83 joined paths and the selected backend lockfile lacked registry artifact identifiers for its corroborated paths. The second yielded one joined root path and 725 joined paths across its selected workspaces. Under the prespecified completeness rule, both root-to-importer coverage changes remained unresolved. A post hoc check found that the second root had exactly one mandatory dependency, no mandatory children in the lock or exact-version metadata, and consistent artifact and maintainership records. With this confirmed root zero, selecting the recipe's importers produces a documented gain for the name portfolio and all three version alternatives. This establishes one importer-sensitive case under the declared dependency scope. Six named candidates remain unresolved, and the 99-recipe sample supplies no population-wide effect estimate.

\subsection{Release-delivered dependency evidence}
We followed the 39 recipes in the fresh sample that declared a release asset and had no matched JavaScript package command. Literal setup instructions identify 38 GitHub assets across 37 recipes, using amd64 resource fields where specified. The release API identifies 37 assets. Thirty reported SHA-256 values agree with the recipe pins, six lack an API digest, and one differs. Public attestation lookups for the 30 agreeing digests return subject-matching statements for six assets. Three SLSA provenance statements declare source commits for their release repositories. We parse the statements without verifying signatures. A separate CycloneDX document identifies the selected Java runtime by filename and exact digest. That inventory contains no npm package URLs.

We downloaded the 15 identified assets within a 12\,MiB cap. Fourteen matched their recipe checksum, including five with no API digest. The mismatching download agreed with the API digest. Eleven matching ZIP or gzip TAR archives could be inspected. One contained 198 bundled npm manifests whose metadata lacked the registry artifact identifiers required for the metadata join. Another contained an application manifest and a Yarn Classic lockfile. These two releases support separate checks on declared dependencies and on packaged files.

\subsection{Declared dependencies in a pinned release}
The Pterodactyl panel archive is the one inspected artifact with both a Node manifest and a Yarn Classic lockfile. Its checksum agrees with the recipe, and both dependency files match the commit resolved from its source release tag byte for byte. We processed both views with the same reader and registry snapshot. Of 252 exact-version requests, all returned usable maintainer lists. The path checks corroborate 224 package versions in each view and document hits under the name portfolio and all three version alternatives. The equal results reflect the identical metadata inputs. The archive also includes 17 JavaScript files under its public assets directory. The lockfile evidence does not trace which of its declared dependencies are present in those files.

\subsection{Evidence in shipped JavaScript and dependency files}
We statically parsed all 17 public JavaScript assets in the Pterodactyl archive. The assets directory contains no source-map files, and the JavaScript contains no source-map directive comments. The sole \texttt{sourceMappingURL} occurrence is inside a string literal. The tagged Webpack configuration disables source maps by default in production, with an environment setting that can override the default. The archive and configuration thus delimit the available evidence for tracing bundled modules in this release.

The Mongo Express archive preserves a \texttt{node\_modules} tree. The packaged versions of all 39 regular direct dependencies satisfy the application's declared ranges. For each of these versions, we checked the npm tarball against the registry's SHA-512 integrity value. We then compared the relative paths and SHA-256 hashes of regular files within each dependency directory, excluding nested \texttt{node\_modules} in both archives. All 39 dependencies matched across 1,571 file paths. Using maintainer lists observed for those versions, the name portfolio and all three version alternatives reached the application through a dependency whose files were present in the release. This diagnostic measures file inclusion for direct dependencies in one archive. Runtime loading and transitive completeness remain outside its scope.
